% Minimal retained preamble for S3 SIGMA 2022 article 023 cards taken
% from sources/sigma-2022-023-post-lie-magnus-expansion-and-bch-recursion/
% pinned-original-source.tex (doi:10.3842/SIGMA.2022.023, SIGMA 18 (2022),
% CC BY 4.0).
%
% Omitted and why: the publisher documentclass{sigma} block (source
% lines 1-7 with its ifpdf branch) - the sigma class is unavailable in
% the toolchain and the assembler supplies the article class
% (unavailable-class precedent). The \usepackage{wasysym} line (source
% line 9): no delivered span of any card of this batch uses a wasysym-only
% symbol (the \rhd of the delivered spans is the amssymb relation, and
% the author's wasysym-dependent \rightarrowRHD/\longrightarrowRHD
% arrows, source lines 13-14, are unused by every delivered span).
% \lbutcher (source line 12) is unused by every delivered span. Of the
% author macro block only the declarations used by the delivered spans
% are copied verbatim below, in source order: the \nc shorthand (source
% line 11, needed because the copied macro lines are written with it),
% \Cal (line 41), \mop (line 42), \wt (line 46; the delivered spans use
% the LaTeX-standard \widehat/\widetilde directly, so the author's
% \wh shorthand at line 47 is not needed). Omitted as unused by every
% delivered span: \vep/\ep (29), \sigmat (30), \ostar (31), \grad (33),
% \fil (34), the blackboard-letter block \BA-\ZZ (36-40), \mopl (43),
% \frakg (44), \g (45), \wh (47), \un (48), \fleche (50) and any
% further macro-block lines. The author equation numbering (source line
% 16) and the author theorem declarations (source lines 18-27, master
% counter Theorem sectionwise with Corollary/Lemma/Proposition hanging on
% it, the starred Theorem* at line 19, and the definition-style group
% with its brace pair on lines 23 and 27) are copied verbatim below. No
% per-card \usepackage additions are needed (CONTROL 8.8.162: nothing
% outside the amsmath/amssymb/amsthm baseline occurs in the delivered
% spans).
\usepackage{amsmath,amssymb,amsthm}
\numberwithin{equation}{section}

\newtheorem{Theorem}{Theorem}[section]
\newtheorem*{Theorem*}{Theorem}
\newtheorem{Corollary}[Theorem]{Corollary}
\newtheorem{Lemma}[Theorem]{Lemma}
\newtheorem{Proposition}[Theorem]{Proposition}
 { \theoremstyle{definition}
\newtheorem{Definition}[Theorem]{Definition}
\newtheorem{Note}[Theorem]{Note}
\newtheorem{Example}[Theorem]{Example}
\newtheorem{Remark}[Theorem]{Remark} }

\newcommand{\nc}{\newcommand}
\nc{\Cal}[1]{{\mathcal {#1}}}
\nc{\mop}[1]{\mathop{\hbox {\rm #1} }}
\nc{\wt}{\widetilde}
